\chapter{Dialectical Behaviour Therapy}

\section*{Definition and Overview}
Dialectical behaviour therapy (DBT) is a structured psychological therapy developed to help people with intense emotional instability, particularly those with borderline personality disorder. It combines cognitive-behavioural techniques with mindfulness and acceptance, balancing change with validation. DBT teaches practical skills in four areas: mindfulness, distress tolerance, emotion regulation, and interpersonal effectiveness. It is delivered in individual therapy and skills groups, and has strong evidence for reducing self-harm and suicidal behaviour. This chapter explains how DBT works, who it helps, and what treatment involves.

\section*{Epidemiology}
DBT was developed for borderline personality disorder, which affects around 1--2\% of the population and is associated with repeated self-harm, suicidal crises, and high use of mental health services. DBT has also been adapted for other conditions involving emotional dysregulation, including eating disorders, substance use, and treatment-resistant depression. Access to DBT in Australia has grown through public mental health services and private practice, although waiting lists can be long.

\section*{Aetiology and Pathophysiology}
\begin{itemize}
    \item \textbf{Emotional vulnerability:} biological sensitivity to emotions, with intense reactions and slow return to baseline
    \item \textbf{Invalidating environment:} a background in which emotions are dismissed or punished, failing to teach emotion regulation
    \item \textbf{Transactional model:} the interaction of biological vulnerability and an invalidating environment produces emotional dysregulation
    \item \textbf{Skill deficits:} difficulty tolerating distress, regulating emotions, and managing relationships
    \item \textbf{Behavioural patterns:} self-harm and impulsive behaviour are often reinforced by relief from distress, and DBT targets these patterns
\end{itemize}

\section*{Clinical Features}
\begin{itemize}
    \item Intense, rapidly shifting emotions that feel overwhelming
    \item Recurrent suicidal thoughts, self-harm, or impulsive behaviour
    \item Unstable relationships with idealisation and devaluation
    \item Chronic feelings of emptiness and fear of abandonment
    \item Difficulty tolerating distress, with crises triggered by minor events
    \item Impulsivity in spending, substance use, eating, or driving
\end{itemize}

\section*{Differential Diagnosis}
\begin{itemize}
    \item Bipolar disorder, in which mood swings are episodic and less reactive to events
    \item Major depressive and anxiety disorders
    \item Post-traumatic stress disorder, including complex PTSD
    \item Attention deficit hyperactivity disorder, which can mimic impulsivity
    \item Substance use disorders, which can cause emotional instability
\end{itemize}

\section*{When to Seek Medical Care}
Anyone experiencing intense emotional distress, repeated crises, or thoughts of self-harm should seek help from a GP, psychologist, or mental health service. A GP can arrange referral to a psychologist trained in DBT.

\textbf{Call 000 (or local emergency services) for:}
\begin{itemize}
    \item Immediate risk of self-harm or suicide
    \item Any crisis in which the person is in danger or cannot keep themselves safe
\end{itemize}
\textbf{Support lines:} Lifeline (13 11 14) and Suicide Call Back Service (1300 659 467) are available 24 hours.

\section*{Diagnosis and Investigations}
\begin{itemize}
    \item \textbf{Structured clinical interview:} assessment of mood, behaviour, relationships, and history of self-harm
    \item \textbf{Diagnostic assessment:} confirmation of borderline personality disorder or another indication for DBT
    \item \textbf{Risk assessment:} careful evaluation of suicide risk and safety planning
    \item \textbf{Baseline measures:} questionnaires tracking emotional dysregulation, self-harm, and quality of life
    \item \textbf{Medical review:} to exclude physical causes of mood disturbance
\end{itemize}

\section*{Management}
\begin{itemize}
    \item \textbf{Individual therapy:} weekly sessions addressing motivation, skills application, and obstacles
    \item \textbf{Skills group:} weekly group teaching of mindfulness, distress tolerance, emotion regulation, and interpersonal effectiveness
    \item \textbf{Phone coaching:} between-session telephone support to apply skills in real-time crises
    \item \textbf{Therapist consultation team:} therapists meet regularly to maintain quality and prevent burnout
    \item \textbf{Commitment strategies:} explicit agreements to work toward reducing life-threatening behaviours
    \item \textbf{Adjunctive treatment:} medication for co-occurring depression, anxiety, or insomnia, and treatment for substance use where present
    \item \textbf{Adapted DBT:} modified programmes for adolescents and specific conditions
\end{itemize}

\section*{Complications}
\begin{itemize}
    \item Ongoing risk of self-harm and suicide without treatment
    \item Drop-out from therapy, which can undermine progress
    \item Deterioration of relationships, work, and physical health from instability
    \item High use of emergency services during crises
    \item Co-occurring conditions complicating recovery
\end{itemize}

\section*{Prognosis and Outcomes}
DBT is one of the best-evidenced treatments for borderline personality disorder. Trials show significant reductions in self-harm, suicidal behaviour, hospital admissions, and emergency presentations. With sustained participation, usually a year or more, most people develop lasting skills in managing emotions and relationships. Long-term studies show continued improvement over years, and many people go on to live stable, fulfilling lives.

\section*{Prevention and Self-Care}
\begin{itemize}
    \item Practise mindfulness to observe emotions without being overwhelmed
    \item Build a "crisis survival kit" of distress-tolerance skills
    \item Establish a regular routine of sleep, meals, and activity
    \item Maintain a safety plan with trusted contacts and crisis numbers
    \item Communicate needs clearly and practise assertiveness
    \item Attend appointments consistently and complete skills homework
    \item Seek help early when distress rises
\end{itemize}

\section*{Sources and Further Reading}
The following reputable sources provide further information for healthcare students and patients seeking reliable, current advice about dialectical behaviour therapy.

\begin{itemize}
    \item Australian Psychological Society. \textit{Dialectical behaviour therapy: an overview.} https://psychology.org.au/
    \item NHS. \textit{Treatment for borderline personality disorder.} https://www.nhs.uk/mental-health/conditions/borderline-personality-disorder/
    \item National Institute of Mental Health. \textit{Borderline personality disorder and treatment.} https://www.nimh.nih.gov/
    \item Psychology Today. \textit{Dialectical behaviour therapy: what it is and how it works.} https://www.psychologytoday.com/
    \item MSD Manual. \textit{Borderline personality disorder: professional overview.} https://www.msdmanuals.com/
    \item Beyond Blue. \textit{Support for difficult emotions and mental health.} https://www.beyondblue.org.au/
\end{itemize}
